% Session 3 -- the SSH invariant made visible: the winding of d(k)=t+e^{ik}.
% Placement: sec. "The SSH model", after equation (3.6) (the Zak phase).
%
% LAYOUT NOTE: every annotation lives inside its own axis and is placed in
% "axis cs:" coordinates. The complex-plane panels use "scale only axis" with
% identical x- and y-ranges so that the circles come out round.
\begin{aqmfigure}{fig:s03:winding}{A drawing of the one-dimensional invariant. The top row shows the
off-diagonal element $d(k)=t+e^{ik}$ of the SSH Bloch Hamiltonian (3.5) in the complex plane, as $k$
goes once through the Brillouin zone, and the arrow shows the direction of increasing $k$. It is
always the unit circle, but centred at $(t,0)$, so it encloses the origin for $t<1$ and does not
enclose it for $t>1$. The winding number of this curve around the origin is exactly $\gamma/\pi$,
which is the Zak phase (3.6) in units of $\pi$. The bottom row shows the corresponding bands
$E_\pm(k)=\pm|d(k)|$, with gap $2|1-t|$. The two rows give the same information in two ways. The only
way to unwind the curve is to move it across the origin, and $d(k)=0$ is where the gap closes, so the
invariant cannot change while the system stays insulating. Note that (a) and (c) have the same gap
and look like ordinary spectra. The difference between them cannot be seen in the energies and is
only in the wavefunctions, which is the main point of (3.4). By the bulk-boundary correspondence,
$\gamma=\pi$ is also exactly the condition for the zero-energy end modes of
figure~\ref{fig:s03:ssh}.}

\begin{tikzpicture}

% =====================================================================
% top row: d(k) in the complex plane. Columns at x = 0, 5.3, 10.6 cm.
% =====================================================================

% ---------- (a) t = 0.5 : winding 1 ----------------------------------
\begin{axis}[
  aqmplot, scale only axis, width=3.5cm, height=3.5cm,
  at={(0cm,0cm)}, anchor=north west,
  axis lines=middle,
  xmin=-1.45, xmax=2.85, ymin=-2.15, ymax=2.15,
  xtick={1,2}, ytick={1},
  xlabel={$\mathrm{Re}\,d$}, ylabel={$\mathrm{Im}\,d$},
  xlabel style={at={(axis description cs:1.03,0.5)}, anchor=west, font=\footnotesize},
  ylabel style={font=\footnotesize},
  title={(a) $t=0.5$},
]
\addplot[aqmcurve, domain=0:360, samples=90] ({0.5+cos(x)},{sin(x)});
\draw[aqm, draw=aqmaccent, -{Stealth[length=5pt,width=4pt]}]
  (axis cs:0.72,0.97) -- (axis cs:0.24,0.97);
\draw[aqm, draw=aqmbroken, line width=1.2pt] (axis cs:-0.14,-0.14) -- (axis cs:0.14,0.14);
\draw[aqm, draw=aqmbroken, line width=1.2pt] (axis cs:-0.14,0.14) -- (axis cs:0.14,-0.14);
\node[aqmlabsmall, text=aqmbroken, anchor=north east] at (axis cs:-0.10,-0.14) {$0$};
\node[aqmlabsmall, anchor=north, align=center] at (axis cs:0.7,-2.32)
  {origin \textbf{enclosed}\\[-1pt]winding $=1$, $\gamma=\pi$\\[-1pt]\textbf{topological}};
\end{axis}

% ---------- (b) t = 1 : critical --------------------------------------
\begin{axis}[
  aqmplot, scale only axis, width=3.5cm, height=3.5cm,
  at={(5.3cm,0cm)}, anchor=north west,
  axis lines=middle,
  xmin=-1.45, xmax=2.85, ymin=-2.15, ymax=2.15,
  xtick={1}, ytick={1},
  xlabel={$\mathrm{Re}\,d$},
  xlabel style={at={(axis description cs:1.03,0.5)}, anchor=west, font=\footnotesize},
  title={(b) $t=1$},
]
\addplot[aqmcurvebroken, domain=0:360, samples=90] ({1+cos(x)},{sin(x)});
\draw[aqm, draw=aqmbroken, -{Stealth[length=5pt,width=4pt]}]
  (axis cs:1.22,0.97) -- (axis cs:0.74,0.97);
\draw[aqm, draw=aqmbroken, line width=1.2pt] (axis cs:-0.14,-0.14) -- (axis cs:0.14,0.14);
\draw[aqm, draw=aqmbroken, line width=1.2pt] (axis cs:-0.14,0.14) -- (axis cs:0.14,-0.14);
\node[aqmlabsmall, text=aqmbroken, anchor=north east] at (axis cs:-0.10,-0.14) {$0$};
\node[aqmlabsmall, anchor=north, align=center, text=aqmbroken] at (axis cs:0.9,-2.32)
  {curve \textbf{hits} the origin\\[-1pt]winding undefined\\[-1pt]\textbf{transition}};
\end{axis}

% ---------- (c) t = 1.5 : winding 0 -----------------------------------
\begin{axis}[
  aqmplot, scale only axis, width=3.5cm, height=3.5cm,
  at={(10.6cm,0cm)}, anchor=north west,
  axis lines=middle,
  xmin=-1.45, xmax=2.85, ymin=-2.15, ymax=2.15,
  xtick={1,2}, ytick={1},
  xlabel={$\mathrm{Re}\,d$},
  xlabel style={at={(axis description cs:1.03,0.5)}, anchor=west, font=\footnotesize},
  title={(c) $t=1.5$},
]
\addplot[aqmcurve, domain=0:360, samples=90] ({1.5+cos(x)},{sin(x)});
\draw[aqm, draw=aqmaccent, -{Stealth[length=5pt,width=4pt]}]
  (axis cs:1.72,0.97) -- (axis cs:1.24,0.97);
\draw[aqm, draw=aqmmuted, line width=1.2pt] (axis cs:-0.14,-0.14) -- (axis cs:0.14,0.14);
\draw[aqm, draw=aqmmuted, line width=1.2pt] (axis cs:-0.14,0.14) -- (axis cs:0.14,-0.14);
\node[aqmlabsmall, text=aqmmuted, anchor=north east] at (axis cs:-0.10,-0.14) {$0$};
\node[aqmlabsmall, anchor=north, align=center] at (axis cs:0.9,-2.32)
  {origin \textbf{outside}\\[-1pt]winding $=0$, $\gamma=0$\\[-1pt]\textbf{trivial}};
\end{axis}

% =====================================================================
% bottom row: the bands E = +- |d(k)|
% =====================================================================

\begin{axis}[
  aqmplot, scale only axis, width=3.5cm, height=2.6cm,
  at={(0cm,-6.1cm)}, anchor=north west,
  axis lines=middle,
  xmin=-3.9, xmax=3.9, ymin=-3.1, ymax=3.1,
  xtick={-3.1416,3.1416}, xticklabels={$-\pi$,$\pi$},
  ytick=\empty,
  xlabel={$k$}, ylabel={$E$},
  xlabel style={at={(axis description cs:1.03,0.5)}, anchor=west, font=\footnotesize},
  ylabel style={font=\footnotesize},
]
\addplot[aqmcurve, domain=-3.1416:3.1416, samples=160]  {sqrt(1.25+cos(deg(x)))};
\addplot[aqmcurve, domain=-3.1416:3.1416, samples=160] {-sqrt(1.25+cos(deg(x)))};
\node[aqmlabsmall, anchor=north, align=center] at (axis cs:0,-3.35) {gap $=2|1-t|=1$};
\end{axis}

\begin{axis}[
  aqmplot, scale only axis, width=3.5cm, height=2.6cm,
  at={(5.3cm,-6.1cm)}, anchor=north west,
  axis lines=middle,
  xmin=-3.9, xmax=3.9, ymin=-3.1, ymax=3.1,
  xtick={-3.1416,3.1416}, xticklabels={$-\pi$,$\pi$},
  ytick=\empty,
  xlabel={$k$},
  xlabel style={at={(axis description cs:1.03,0.5)}, anchor=west, font=\footnotesize},
]
\addplot[aqmcurvebroken, domain=-3.1416:3.1416, samples=200]  {sqrt(2+2*cos(deg(x)))};
\addplot[aqmcurvebroken, domain=-3.1416:3.1416, samples=200] {-sqrt(2+2*cos(deg(x)))};
\node[circle, draw=aqmbroken, fill=white, line width=\aqmlinewidth,
      minimum size=6pt, inner sep=0pt] at (axis cs:3.1416,0) {};
\node[aqmlabsmall, anchor=north, align=center, text=aqmbroken] at (axis cs:0,-3.35)
  {gap closes at $k=\pi$};
\end{axis}

\begin{axis}[
  aqmplot, scale only axis, width=3.5cm, height=2.6cm,
  at={(10.6cm,-6.1cm)}, anchor=north west,
  axis lines=middle,
  xmin=-3.9, xmax=3.9, ymin=-3.1, ymax=3.1,
  xtick={-3.1416,3.1416}, xticklabels={$-\pi$,$\pi$},
  ytick=\empty,
  xlabel={$k$},
  xlabel style={at={(axis description cs:1.03,0.5)}, anchor=west, font=\footnotesize},
]
\addplot[aqmcurve, domain=-3.1416:3.1416, samples=160]  {sqrt(3.25+3*cos(deg(x)))};
\addplot[aqmcurve, domain=-3.1416:3.1416, samples=160] {-sqrt(3.25+3*cos(deg(x)))};
\node[aqmlabsmall, anchor=north, align=center] at (axis cs:0,-3.35) {gap $=2|1-t|=1$};
\end{axis}

% ---- the phase-diagram strip tying the three columns together --------
\draw[aqm, draw=black, line width=\aqmthinwidth]
  (0cm,-9.55cm) -- (14.1cm,-9.55cm);
\foreach \x in {1.75,7.05,12.35} {
  \draw[aqm, draw=black, line width=\aqmthinwidth] (\x cm,-9.46cm) -- (\x cm,-9.64cm);
}
\node[aqmlabsmall, anchor=north] at (1.75cm,-9.70cm) {$t<1$: topological};
\node[aqmlabsmall, anchor=north, text=aqmbroken] at (7.05cm,-9.70cm) {$t=1$: gap closes};
\node[aqmlabsmall, anchor=north] at (12.35cm,-9.70cm) {$t>1$: trivial};

\end{tikzpicture}

\end{aqmfigure}
